\documentclass[10pt,a4paper]{article}

\usepackage[margin=2.5cm]{geometry}
\usepackage{graphicx}
\usepackage{booktabs}
\usepackage{amsmath, amssymb}
\usepackage{hyperref}
\usepackage{cleveref}
\usepackage{natbib}
\usepackage{caption}

\title{When Do LSTMs Prefer Spurious Signals?\\
       A Study of Trigger Strength and Position in Text Classification}
\author{Author 1, Author 2, Author 3 \\
        COMP3242 Deep Learning, Semester 1 2026}
\date{June 2026}

\begin{document}
\maketitle

\begin{abstract}
We study shortcut learning in sequence models trained for sentiment
classification on IMDb. By injecting a controlled artificial trigger token
into a fraction of one class during training, we measure the rate at which
LSTM and Transformer models adopt the spurious signal as a function of its
frequency and position in the input sequence. We find that\ldots
\textit{[fill in after experiments].}
\end{abstract}

\section{Introduction}
\label{sec:intro}

Shortcut learning---the tendency of neural networks to rely on spurious
correlations rather than the intended task signal---is a well-documented
failure mode \citep{geirhos2020}. \textit{[Person 3 to expand: motivate why
this matters; introduce IMDb sentiment as the testbed; state research
questions and hypotheses.]}

\paragraph{Contributions.}
\begin{itemize}
  \item A controlled study of trigger-strength sensitivity in LSTM and
        Transformer text classifiers.
  \item A direct comparison of position sensitivity between architectures
        with and without an explicit positional inductive bias.
  \item Empirical evidence for \textit{[insert main finding]}.
\end{itemize}

\section{Method}
\label{sec:method}

\subsection{Dataset and preprocessing}
\textit{[Person 1 to write: IMDb dataset, vocab construction, max sequence
length, train/val/test split.]}

\subsection{Trigger injection}
We inject an artificial token \texttt{qzx} into a fraction $p$ of training
examples whose label is positive. The token is reserved in the vocabulary so
it is not mapped to \texttt{<unk>}. \textit{[Person 1 to expand: position
semantics, edge-case handling, why we inject only into one class.]}

\subsection{Models}
\textit{[Person 2 to write: LSTM architecture (1 layer, hidden=128); Transformer
(2 layers, 4 heads, $d=128$, sinusoidal positional encoding). Justify model
sizes by CPU compute budget.]}

\subsection{Evaluation protocol}
We evaluate every trained model under four conditions: (i) the unmodified test
set, (ii) the test set with all occurrences of the trigger token removed,
(iii) the test set with the trigger injected into every example, and (iv) a
flip-rate measurement on the negative class only. \textit{[Person 3 to
formalise each metric.]}

\section{Experiments}
\label{sec:experiments}

We study three axes: trigger strength $p \in \{0, 0.25, 0.5, 0.75\}$, trigger
position $\in \{\text{start}, \text{middle}, \text{end}\}$, and architecture
$\in \{\text{LSTM}, \text{Transformer}\}$. Position experiments are run only
at $p=0.5$ (the diagnostic slice). All configurations are trained with three
random seeds; we report mean $\pm$ standard deviation throughout.

\section{Results}
\label{sec:results}

\subsection{Trigger strength}
\textit{[Insert Figure~\ref{fig:strength}: accuracy vs.\ trigger strength,
two lines (with-trigger test, no-trigger test), separate panels for LSTM
and Transformer. Discuss the phase-transition point.]}

\subsection{Position sensitivity}
\textit{[Insert Figure~\ref{fig:position}: bar chart of flip rate by position,
grouped by architecture. Discuss recency bias.]}

\subsection{Flip rate}
\textit{[Insert Figure~\ref{fig:flip}.]}

\section{Discussion}
\label{sec:discussion}

\textit{[Mechanistic interpretation: why does the LSTM behave this way?
Connect to gating dynamics. Why is the Transformer position-invariant?
Sinusoidal PE provides position information uniformly across the sequence.]}

\paragraph{Limitations.}
\textit{[Discuss: single dataset, single artificial trigger token, CPU compute
budget restricting model size, did not vary trigger token form.]}

\paragraph{Future work.}
\textit{[Natural-word triggers, mechanistic analysis of attention weights and
LSTM gates, comparison to learned positional encodings.]}

\section{Conclusion}
\label{sec:conclusion}

\textit{[One paragraph restating the headline finding.]}

\section*{Author contributions}

\textbf{Author 1} (uXXXXXXX): data pipeline, trigger injection, configuration
system. Wrote the Method section.

\textbf{Author 2} (uXXXXXXX): LSTM and Transformer implementations, training
infrastructure, sweep orchestration. Wrote the Experiments section.

\textbf{Author 3} (uXXXXXXX): evaluation harness, plotting, analysis.
Wrote the Results and Discussion sections.

All authors contributed to the Introduction, Conclusion, and final editing.

\section*{AI assistance declaration}

\textit{[Itemise any use of generative AI tools per COMP3242 guidelines:
which tool, which sections, what kind of help (planning, code review,
prose editing, etc.).]}

\bibliographystyle{plainnat}
\bibliography{references}

\end{document}
